% ..............................................................................
% Slide 9 -- H3: are the learned features stable and distinct per neuron?
% Source content: contents/slide-09-h3.md  (bullets are FINAL -- render as-is)
% ------------------------------------------------------------------------------
\section{H3: Feature quality}

\begin{frame}[t]{H3 --- Are the learned features stable and distinct per neuron?}
% ---- Figure: full width across the top -----------------------------------
\centering
\includegraphics[width=\linewidth,height=0.42\textheight,keepaspectratio]%
{contents/images/h3-stability.png}

\vfill
% ---- Bullets: lower third, two compact columns ---------------------------
{\footnotesize
\begin{columns}[t]
\begin{column}{0.5\textwidth}
\begin{itemize}\itemsep=0.1em
\item \textbf{Hypothesis:} a good representation keeps a neuron looking like
      itself over time, and distinct from other neurons
\item \textbf{Metric (uses no labels):} cosine similarity of the encoder's
      features --- within-neuron (same neuron, different frames) vs
      between-neuron (different neurons)
\end{itemize}
\end{column}
\begin{column}{0.5\textwidth}
\begin{itemize}\itemsep=0.1em
\item \textbf{Gap $=$ within $-$ between}; a large gap means the encoder actually
      tells neurons apart
\item \textbf{Result:} the self-supervised encoder shows a \textbf{clear gap, far
      above random} $\to$ it learned to distinguish neurons \textbf{without ever
      seeing a label} (supervised is strongest)
\end{itemize}
\end{column}
\end{columns}
}
\end{frame}
